\section{Commands and Action Semantics}

\subsection{Well-Formed Commands}

A parsed command is a triple

\[
  c=(action,dns,verbosity).
\]

The well-formedness predicate requires exactly one action and requires a DNS
argument exactly when that action is \code{add}:

\[
\begin{split}
\mathsf{wellFormed}(c)\quad\Longleftrightarrow\quad
  &c.action\in\Action\\
  &{}\land\bigl(c.action=\mathsf{add}
    \Longleftrightarrow c.dns=\Some{d}\land\Valid(d)\bigr).
\end{split}
\]

Verbosity is a modifier, not an action.  Short and long forms are aliases:

\begin{align*}
\Parse(\text{``-a''},d) &= \Parse(\text{``--add''},d),\\
\Parse(\text{``-f''}) &= \Parse(\text{``--flush''}),\\
\Parse(\text{``-h''}) &= \Parse(\text{``--help''}),\\
\Parse(\text{``-v''},a) &= \Parse(\text{``--verbose''},a).
\end{align*}

Any command containing two actions is rejected, for example

\[
  \Parse(\text{``--flush prep''})=\Failure.
\]

\subsection{Add}

For a valid DNS name $d$, the successful abstract effect is

\[
\Add{d}(S)=S[
  A\coloneqq\Insert{d}(S.A),
  B\coloneqq\Some{S.H},
  H\coloneqq\Remove{d}(S.H)].
\]

The active hosts file is backed up immediately before replacement, and its
metadata is preserved:

\[
  \Add{d}(S).M=S.M.
\]

An invalid name has no intended state transition:

\[
  \neg\Valid(d)\Longrightarrow
  \Exec(\mathsf{add}(d),S)=(\Failure,S,t)
\]

for some diagnostic trace $t$.

\subsection{Prep}

Let $\HBlock(A,E)$ denote the hosts contents produced by the external
\code{hblock} program from allow list $A$ and external environment $E$.
Subject to successful confirmation when $O$ already exists,

\[
  \Prep_E(S)=S[O\coloneqq\Some{S.H},
                 H\coloneqq\HBlock(S.A,E)].
\]

The environment parameter represents block sources, network results, hblock
version, and configuration outside this repository.

\subsection{Restore}

Restore is defined only when the original backup exists:

\[
\Restore(S)=
  \begin{cases}
    S[H\coloneqq H_O], & S.O=\Some{H_O},\\
    \Failure, & S.O=\None.
  \end{cases}
\]

Existing active contents require confirmation before replacement.

\subsection{Flush}

Let $\Fresh$ denote the desired cache state with a running
\code{mDNSResponder}:

\[
\Flush(S)=
  \begin{cases}
    S[K\coloneqq\Fresh], & \text{on macOS and command success},\\
    \Failure, & \text{on other platforms or command failure}.
  \end{cases}
\]

Flush does not intentionally modify $H$, $O$, $A$, $B$, or $M$.
